% ============================================================================
% WORKBOOK — Section 6.2: Reproducing Scale Drawings at a Different Scale
% CCSS 7.G.A.1
% Practice-focused reinforcement for the study guide topic
% ============================================================================
\section{Reproducing Scale Drawings at a Different Scale}
\topicTitle{Reproducing Scale Drawings at a Different Scale}

\begin{quickReview}{Changing the Scale}
To redraw a scale drawing at a \textbf{new scale}, find how the two scales relate:

$$\text{new drawing length} = \text{old drawing length} \times \frac{\text{old scale factor}}{\text{new scale factor}}$$

\begin{itemize}[leftmargin=*]
    \item A \textbf{larger scale} (fewer real units per cm) makes the drawing \textbf{bigger}.
    \item A \textbf{smaller scale} (more real units per cm) makes the drawing \textbf{smaller}.
\end{itemize}

\medskip
\textbf{Example:} Old scale: $1$ cm $= 10$ m. New scale: $1$ cm $= 5$ m. Scale ratio: $\frac{10}{5} = 2$. A $3$ cm line becomes $3 \times 2 = 6$ cm on the new drawing.
\end{quickReview}

% ── Warm-Up ──────────────────────────────────────────────────────────────────
\begin{practiceBox}[funGreen]{\faSun~Warm-Up}

\practiceHeader[funGreen]{Find the Scale Ratio}
Divide to find how the two scales relate.

\begin{multicols}{2}
\prob Old: $1$ cm $= 20$ m. New: $1$ cm $= 10$ m. Scale ratio? \answerBlank[2cm]
\answerExplain{$2$}{Divide: $\frac{20}{10} = 2$. Every old measurement is multiplied by $2$ on the new drawing.}

\prob Old: $1$ cm $= 50$ km. New: $1$ cm $= 25$ km. Scale ratio? \answerBlank[2cm]
\answerExplain{$2$}{Divide: $\frac{50}{25} = 2$. The new drawing is twice as large.}

\prob Old: $1$ cm $= 4$ m. New: $1$ cm $= 12$ m. Scale ratio? \answerBlank[2cm]
\answerExplain{$\frac{1}{3}$}{Divide: $\frac{4}{12} = \frac{1}{3}$. The new drawing is one-third as large.}

\prob Old: $1$ in $= 6$ ft. New: $1$ in $= 3$ ft. Scale ratio? \answerBlank[2cm]
\answerExplain{$2$}{Divide: $\frac{6}{3} = 2$. Every old measurement is doubled.}

\prob Old: $1$ cm $= 100$ m. New: $1$ cm $= 50$ m. Scale ratio? \answerBlank[2cm]
\answerExplain{$2$}{Divide: $\frac{100}{50} = 2$. The new drawing is twice as large.}

\prob Old: $1$ in $= 8$ ft. New: $1$ in $= 24$ ft. Scale ratio? \answerBlank[2cm]
\answerExplain{$\frac{1}{3}$}{Divide: $\frac{8}{24} = \frac{1}{3}$. The new drawing is one-third as large because the new scale covers more real distance per inch.}
\end{multicols}
\end{practiceBox}

% ── Practice A: Rescaling Lengths ────────────────────────────────────────────
\begin{practiceBox}{Rescaling Lengths}

\practiceHeader{Use the scale ratio to find the new drawing length.}

\prob Old scale: $1$ cm $= 5$ m. New scale: $1$ cm $= 10$ m. Old length: $8$ cm. New length? \answerBlank[3cm]
\answerExplain{$4$ cm}{Scale ratio: $\frac{5}{10} = \frac{1}{2}$. New length: $8 \times \frac{1}{2} = 4$ cm. The new scale covers more real distance per cm, so the drawing shrinks.}

\prob Old scale: $1$ cm $= 30$ km. New scale: $1$ cm $= 15$ km. Old length: $6$ cm. New length? \answerBlank[3cm]
\answerExplain{$12$ cm}{Scale ratio: $\frac{30}{15} = 2$. New length: $6 \times 2 = 12$ cm.}

\prob Old scale: $1$ in $= 4$ ft. New scale: $1$ in $= 2$ ft. Old length: $5$ in. New length? \answerBlank[3cm]
\answerExplain{$10$ in}{Scale ratio: $\frac{4}{2} = 2$. New length: $5 \times 2 = 10$ in.}

\prob Old scale: $1$ cm $= 40$ m. New scale: $1$ cm $= 8$ m. Old length: $3$ cm. New length? \answerBlank[3cm]
\answerExplain{$15$ cm}{Scale ratio: $\frac{40}{8} = 5$. New length: $3 \times 5 = 15$ cm.}

\prob Old scale: $1:200$. New scale: $1:100$. Old length: $7$ cm. New length? \answerBlank[3cm]
\answerExplain{$14$ cm}{Scale ratio: $\frac{200}{100} = 2$. New length: $7 \times 2 = 14$ cm.}

\prob Old scale: $1:50$. New scale: $1:150$. Old length: $12$ cm. New length? \answerBlank[3cm]
\answerExplain{$4$ cm}{Scale ratio: $\frac{50}{150} = \frac{1}{3}$. New length: $12 \times \frac{1}{3} = 4$ cm.}
\end{practiceBox}

% ── Practice B: Rescaling Rectangles ─────────────────────────────────────────
\begin{practiceBox}[funTeal]{Rescaling Rectangles}

\practiceHeader[funTeal]{Find the new dimensions and the new area on the rescaled drawing.}

\prob A rectangle is $4$ cm by $6$ cm. Old scale: $1$ cm $= 10$ m. New scale: $1$ cm $= 5$ m. \\
New dimensions? \answerBlank[3cm] \quad New drawing area? \answerBlank[3cm]
\answerExplain{$8$ cm by $12$ cm; area $= 96$ cm$^2$}{Scale ratio: $\frac{10}{5} = 2$. New: $4 \times 2 = 8$ cm by $6 \times 2 = 12$ cm. Area: $8 \times 12 = 96$ cm$^2$.}

\prob A rectangle is $9$ cm by $6$ cm. Old scale: $1$ cm $= 3$ m. New scale: $1$ cm $= 9$ m. \\
New dimensions? \answerBlank[3cm] \quad New drawing area? \answerBlank[3cm]
\answerExplain{$3$ cm by $2$ cm; area $= 6$ cm$^2$}{Scale ratio: $\frac{3}{9} = \frac{1}{3}$. New: $9 \times \frac{1}{3} = 3$ cm by $6 \times \frac{1}{3} = 2$ cm. Area: $3 \times 2 = 6$ cm$^2$.}

\prob A square is $5$ cm by $5$ cm. Old scale: $1:20$. New scale: $1:10$. \\
New dimensions? \answerBlank[3cm] \quad New drawing area? \answerBlank[3cm]
\answerExplain{$10$ cm by $10$ cm; area $= 100$ cm$^2$}{Scale ratio: $\frac{20}{10} = 2$. New side: $5 \times 2 = 10$ cm. Area: $10 \times 10 = 100$ cm$^2$.}
\end{practiceBox}

% ── Practice C: True or False ────────────────────────────────────────────────
\begin{practiceBox}[funPurple]{True or False?}

\trueOrFalse{When you redraw a scale drawing at a different scale, the angles in the shapes change.}
\answerExplain{False}{Changing the scale changes only the lengths. All angles remain exactly the same because the shapes stay proportional.}

\trueOrFalse{Going from $1$ cm $= 20$ m to $1$ cm $= 10$ m makes the drawing larger.}
\answerExplain{True}{The new scale represents fewer real meters per cm, so each real-world length needs more cm on the drawing, making the drawing larger.}

\trueOrFalse{If the scale ratio is $3$, areas on the new drawing are $3$ times the old areas.}
\answerExplain{False}{If lengths are multiplied by $3$, areas are multiplied by $3^2 = 9$. Areas always scale by the square of the length scale ratio.}

\trueOrFalse{Going from $1$ cm $= 5$ m to $1$ cm $= 15$ m makes the drawing one-third as large.}
\answerExplain{True}{Scale ratio: $\frac{5}{15} = \frac{1}{3}$. Every length on the new drawing is $\frac{1}{3}$ of the old length, so the drawing shrinks to one-third its original size.}
\end{practiceBox}

% ── Practice D: Word Problems ────────────────────────────────────────────────
\begin{practiceBox}[funBlue]{\faGlobe~Word Problems}

\wordProblem{An architect's floor plan uses the scale $1$ cm $= 2$ m. She needs to present it on a poster at $1$ cm $= 1$ m. A room is $4$ cm by $3$ cm on the original. What are the dimensions on the poster?}{cm by cm}
\answerExplain{$8$ cm by $6$ cm}{Scale ratio: $\frac{2}{1} = 2$. Multiply each dimension: $4 \times 2 = 8$ cm and $3 \times 2 = 6$ cm.}

\wordProblem{A city map uses $1$ cm $= 500$ m. A tourist map uses $1$ cm $= 250$ m. A street is $3$ cm on the city map. How long is it on the tourist map?}{cm}
\answerExplain{$6$ cm}{Scale ratio: $\frac{500}{250} = 2$. New length: $3 \times 2 = 6$ cm. The tourist map is more detailed, so the street appears longer.}

\wordProblem{A model airplane uses the scale $1:72$. A museum wants to display it at the scale $1:36$. The wingspan on the original model is $15$ cm. What is the wingspan on the new model?}{cm}
\answerExplain{$30$ cm}{Scale ratio: $\frac{72}{36} = 2$. New wingspan: $15 \times 2 = 30$ cm.}
\end{practiceBox}

% ── Challenge ────────────────────────────────────────────────────────────────
\begin{challengeBox}
\prob A rectangular park is $6$ cm by $4$ cm on Map A ($1$ cm $= 100$ m). Map B uses $1$ cm $= 200$ m. What is the area of the park on Map B? What is the actual area of the park?
\answerExplain{Map B area: $6$ cm$^2$; Actual area: $240{,}000$ m$^2$}{Scale ratio: $\frac{100}{200} = \frac{1}{2}$. Map B dimensions: $6 \times \frac{1}{2} = 3$ cm by $4 \times \frac{1}{2} = 2$ cm. Map B area: $3 \times 2 = 6$ cm$^2$. Actual dimensions: $6 \times 100 = 600$ m by $4 \times 100 = 400$ m. Actual area: $600 \times 400 = 240{,}000$ m$^2$.}
\end{challengeBox}

\encouragement{Excellent rescaling skills! You can now convert between any two scales with confidence!}
